\capitulo{El examen por dentro}{Qu\'e va a venir, de d\'onde sale cada problema y c\'omo usar esta gu\'ia.}{cap:examen}

\subsection{Lo que revela tu examen del Tema 2}

Tu segundo parcial tuvo \textbf{5 problemas}, uno por subtema, todos de procedimiento. Al compararlo con la Serie 2
se ve de d\'onde sali\'o cada uno:

\begin{center}\small
\begin{tabular}{|>{\centering\arraybackslash}p{0.7cm}|p{6.7cm}|p{7.2cm}|}\hline
\hd{\#} & \hd{Qu\'e pidi\'o} & \hd{De d\'onde sale el molde}\\\hline
1 & Dada $F_X$: media, variancia y una probabilidad & Serie 2, problema 9: misma idea con otra funci\'on\\\hline
2 & Dada $F$ discreta: funci\'on de masa y $P(4<X\le7)$ & \cellcolor{marca}Serie 2, problema 7: \textbf{casi id\'entico}\\\hline
3 & Densidad: media, mediana, moda y sesgo & Serie 2, problemas 13 y 15\\\hline
4 & Hallar $c$, obtener $F_Y$, graficar y usarla & Serie 2, problemas 8 y 13\\\hline
5 & Valor esperado de un costo $120X+10$ & Serie 2, problema 14 y propiedades de $\E$\\\hline
\end{tabular}
\end{center}

\begin{idea}{El patr\'on de la profesora}
\begin{itemize}
\item Toma \textbf{un problema de la serie casi tal cual} y arma los dem\'as con el mismo molde, cambiando la funci\'on o los n\'umeros.
\item Recorre el temario \textbf{en orden}: un problema por subtema.
\item Pide \textbf{incisos cortos y encadenados}: el resultado de a) se usa en b).
\item Califica \textbf{procedimiento y resultado}, y pide que el resultado est\'e resaltado.
\end{itemize}
Conclusi\'on: si dominas los moldes de la Serie 3 y de los ejemplos 3.3 a 3.13 de tus apuntes, no hay sorpresa posible.
\end{idea}

\subsection{Pron\'ostico del examen del Tema 3}

Las probabilidades son \textbf{mi estimaci\'on}, hecha con el patr\'on anterior, el espacio que ocupa cada subtema en tus
apuntes y cu\'antas veces aparece en la serie. No son un dato oficial.

\begin{center}
\begin{tikzpicture}[x=0.072cm,y=-0.66cm]
\foreach \lab/\p [count=\k from 0] in {
  {Tabla discreta: eventos, marginales y condicional}/95,
  {Densidad: hallar $k$ y una probabilidad sobre una regi\'on}/95,
  {Marginales y condicionales continuas}/85,
  {Covariancia y coeficiente de correlaci\'on}/85,
  {?`Son independientes? (como inciso)}/75,
  {$\E$ y $\Var$ de una combinaci\'on lineal}/70,
  {Valor esperado con contexto (costo, sobrante)}/65,
  {Curva de regresi\'on}/45,
  {Construir la tabla desde un experimento}/30,
  {Funci\'on de distribuci\'on conjunta}/25,
  {Generadora de momentos de una suma}/10}{
  \node[anchor=east,font=\small,text=tinta] at (-2,\k) {\lab};
  \fill[tinta!9] (0,\k-0.27) rectangle (100,\k+0.27);
  \ifnum\p>79 \fill[colY] (0,\k-0.27) rectangle (\p,\k+0.27);
  \else\ifnum\p>59 \fill[colY!65] (0,\k-0.27) rectangle (\p,\k+0.27);
  \else \fill[tinta!45] (0,\k-0.27) rectangle (\p,\k+0.27);\fi\fi
  \node[anchor=west,font=\sffamily\bfseries\small,text=tinta] at (\p+1,\k) {\p\,\%};
}
\end{tikzpicture}
\end{center}

\textbf{El examen m\'as probable} tiene estos cinco problemas:

\begin{enumerate}[label=\protect\circnum{\arabic*},leftmargin=10mm,labelsep=3mm]
\item Una \textbf{tabla} conjunta: probabilidad de un evento, marginales, una condicional y si son independientes.
\item Una \textbf{densidad} con constante por determinar y una probabilidad sobre un tri\'angulo.
\item La \textbf{marginal y la condicional} de una densidad, con una probabilidad condicional.
\item \textbf{Covariancia y correlaci\'on} (de una tabla o de una densidad).
\item Una \textbf{combinaci\'on lineal} o un \textbf{valor esperado con contexto}.
\end{enumerate}

\subsection{Mapa: del enunciado a lo que tienes que hacer}

Todo el tema son diez acciones. Lee qu\'e pide el inciso, busca la frase en la columna izquierda y haz lo de la derecha.

\begin{center}\small
\rowcolors{2}{papel}{white}
\begin{tabular}{|p{5.6cm}|p{8.1cm}|>{\centering\arraybackslash}p{1.2cm}|}\hline
\hd{Si el inciso dice\ldots} & \hd{\ldots haces esto} & \hd{Cap.}\\\hline
\guillemotleft Probabilidad de que\ldots\guillemotright{} (tabla) & Marcas las celdas que cumplen y las sumas & \ref{cap:discretas}\\\hline
\guillemotleft Probabilidad de que\ldots\guillemotright{} (f\'ormula) & Dibujas la regi\'on y haces una integral doble & \ref{cap:continuas}\\\hline
\guillemotleft Valor de $k$, $c$, $a$\guillemotright & Igualas a 1 la suma o la integral de toda la funci\'on & \ref{cap:continuas}\\\hline
\guillemotleft Funci\'on marginal\guillemotright & Sumas o integras \textbf{la otra} variable & \ref{cap:discretas}, \ref{cap:margcond}\\\hline
\guillemotleft Condicional\guillemotright, \guillemotleft dado que\guillemotright, \guillemotleft si se sabe que\guillemotright & Conjunta entre la marginal de lo que ya conoces & \ref{cap:discretas}, \ref{cap:margcond}\\\hline
\guillemotleft ?`Son independientes?\guillemotright & Comparas la conjunta con el producto de las marginales & \ref{cap:discretas}, \ref{cap:margcond}\\\hline
\guillemotleft Valor esperado\guillemotright, \guillemotleft promedio\guillemotright, \guillemotleft costo esperado\guillemotright & Multiplicas lo que se promedia por la conjunta y sumas o integras & \ref{cap:esperanza}\\\hline
\guillemotleft Curva de regresi\'on\guillemotright, \guillemotleft pron\'ostico\guillemotright & Valor esperado usando la condicional & \ref{cap:esperanza}\\\hline
\guillemotleft Covariancia\guillemotright, \guillemotleft correlaci\'on\guillemotright & Calculas cinco esperanzas y aplicas dos f\'ormulas & \ref{cap:cov}\\\hline
\guillemotleft $\E$ o $\Var$ de $aX+bY+c$\guillemotright & Aplicas las dos f\'ormulas de combinaci\'on lineal & \ref{cap:lineal}\\\hline
\end{tabular}
\rowcolors{1}{}{}
\end{center}

\subsection{C\'omo leer esta gu\'ia}

\begin{center}\small
\begin{tabular}{|p{4.3cm}|p{11.1cm}|}\hline
\hd{Vas a ver} & \hd{Significa}\\\hline
\cx{\textbf{Verde azulado}} y \cy{\textbf{naranja}} & Todo lo que pertenece a $\cx{X}$ va en verde azulado y todo lo de $\cy{Y}$ en naranja, en tablas, f\'ormulas y figuras.\\\hline
\hlt{Marcatextos} & Las celdas de un evento, o la parte de una f\'ormula que cambi\'o en ese paso.\\\hline
\res{\text{Recuadro verde}} & Resultado final de un inciso. En tu examen enci\'erralo igual: las instrucciones lo piden.\\\hline
\circnum{1}\ \circnum{2}\ \circnum{3} & Pasos en orden. Cada paso tiene un t\'itulo que dice \textbf{qu\'e} se hace.\\\hline
\textcolor{colX}{\sffamily\bfseries ?`Por qu\'e?} & La raz\'on del paso, para que lo entiendas y no s\'olo lo copies.\\\hline
\textcolor{verde}{\sffamily\bfseries \ding{51} Verifica} & Una comprobaci\'on de diez segundos que detecta errores.\\\hline
\textcolor{colY}{\sffamily\bfseries TU TURNO} & Ejercicio para hacer sin ver. La respuesta viene con sus n\'umeros intermedios.\\\hline
\end{tabular}
\end{center}

\textbf{Plan de estudio sugerido} (unas 7 horas): cap\'itulos \ref{cap:bases} y \ref{cap:discretas} en la primera sesi\'on;
\ref{cap:continuas} y \ref{cap:margcond} con l\'apiz en la segunda; \ref{cap:esperanza}, \ref{cap:cov} y \ref{cap:lineal} en la tercera;
el simulacro del cap\'itulo \ref{cap:cierre} con reloj el d\'ia anterior al examen.
